\section{未经规划的意图为何经常失败}

\begin{proposition}[未经规划的意图保留了当前优势]
假设在某个决策窗口，有两个未来选项相互竞争：留在当前状态 $s_0$，或尝试转变到 $s_1$。如果没有任何计划去降低 $s_1$ 的障碍项，那么该模型就不会提供新的结构性理由，使得在那个窗口里 $s_1$ 会比现状更容易。
\end{proposition}

\begin{proof}
根据假设，没有任何计划去降低含糊性、提高准备，或以其他方式减少 $s_1$ 的转变负担。因此，模型中与障碍相关的变量并没有朝有利于 $s_1$ 的方向变化。既然当前状态的结构性优势没有被计划削弱，模型也就没有引入新的形式理由，让这次转变相对于现状更容易。
\end{proof}

\begin{corollary}[未经规划的偏好本身并不会逆转惯性]
在这个框架内部，对 $s_1$ 的一种口头偏好，本身并不意味着在下一次决策窗口中，$s_1$ 就会击败当前状态的结构性优势。
\end{corollary}

\begin{proof}
把前一个命题应用到这样一个特例即可：唯一的新元素只是口头认可，而没有发生任何与障碍相关的计划变化。
\end{proof}

\begin{example}[下班后锻炼]
“我下班后应该去锻炼”这句话，还不能算一个可用的计划。它没有解决地点、背包、动作顺序、训练内容、时长和退出条件。在本章的记号里，这些未解决的项目让执行复杂度保持高位，也使目标含糊性几乎没有变化。真正的计划，会在上游把其中一部分变量先解决掉。
\end{example}